\documentclass[10pt,a4paper]{article}

\usepackage[T1]{fontenc}
\usepackage[utf8]{inputenc}
\usepackage{lmodern,amsmath,amssymb,array,tabularx,booktabs,xcolor,enumitem}
\usepackage[margin=20mm]{geometry}
\usepackage[unicode,colorlinks=true,urlcolor=blue,linkcolor=black]{hyperref}

\setlength{\parindent}{0pt}
\setlength{\parskip}{5pt}
\setlength{\emergencystretch}{2em}
\setlength{\tabcolsep}{4pt}
\renewcommand{\arraystretch}{1.10}
\newcolumntype{Y}{>{\raggedright\arraybackslash}X}
\newcommand{\paper}[1]{\href{#1}{link}}

\title{\textbf{Baseline model}\\[2pt]}

\begin{document}
\maketitle
\vspace{-12mm}

%%%%%%%%%%%%%%%%%%%%%%%%%%%%%%%%%%%%%%%%%%%%%%%%%%%%%%
\section{Overview}

This document describes a two-stage stochastic optimization model that decides
(i) which candidate cities a robotaxi operator should enter and
(ii) how many vehicles to deploy in each entered city,
under uncertain robotaxi demand represented by a finite set of scenarios.

The model trades off four cost components:
\begin{itemize}[leftmargin=*,itemsep=1pt]
  \item fixed city-entry costs,
  \item per-vehicle costs,
  \item a penalty for unused service capacity (over-supply), and
  \item a penalty for unserved robotaxi demand, i.e.\ foregone contribution margin (under-supply).
\end{itemize}

\textbf{Evaluation period.} All demand, throughput, and cost quantities are measured
over \emph{the same evaluation period}. The model determines the city/vehicle
deployment to be applied during that period.

%%%%%%%%%%%%%%%%%%%%%%%%%%%%%%%%%%%%%%%%%%%%%%%%%%%%%%
\section{Model}

\subsection{Objective}

\[
\min_{y,q,O,U}\quad
\underbrace{\sum_{c} F_c y_c}_{\text{entry costs}}
+\underbrace{\sum_{c} K_c q_c}_{\text{vehicle costs}}
+\frac{1}{|S|}\sum_{s\in S}\sum_{c\in C}\Big[
\underbrace{c_c^{+}O_c^{s}}_{\text{unused capacity}}
+\underbrace{c_c^{-}U_c^{s}}_{\text{foregone margin}}
\Big].
\]

The superscripts in $c_c^{+}$ and $c_c^{-}$ denote \emph{excess} and \emph{shortage},
not the sign of the cost. Scenarios are weighted equally ($1/|S|$).

\subsection{Constraints}

\[
\begin{aligned}
&0\le q_c\le \bar q_c\, y_c && \forall c\in C,\\
&\textstyle\sum_{c} q_c\le Q, &&\\
&\kappa_c q_c - M_c p_c^{s} = O_c^{s}-U_c^{s} && \forall c\in C,\ s\in S,\\
&y_c\in\{0,1\},\quad q_c\in\mathbb{Z}_{\ge 0}, &&\\
&O_c^{s},\,U_c^{s}\ge 0. &&
\end{aligned}
\]

{\small
\begin{tabularx}{\linewidth}{@{}>{\raggedright\arraybackslash}p{0.32\linewidth}Y@{}}
\toprule
Constraint & Meaning \\ \midrule
$0\le q_c\le\bar q_c y_c$ & No vehicles in a city that is not entered; up to the city-specific limit if it is entered. \\
$\sum_c q_c\le Q$ & Total fleet limit. \\
$\kappa_c q_c-M_c p_c^s=O_c^s-U_c^s$ & Service capacity minus robotaxi demand equals unused capacity minus unmet demand. \\
\bottomrule
\end{tabularx}
}

\subsection{Demand construction}

Each city has a pool of potential ride-hailing usage opportunities $M_c$, which is
split into robotaxi demand and non-robotaxi demand:
\[
D_c^{s}=M_c\,p_c^{s},\qquad D_c^{N,s}=M_c\,(1-p_c^{s}).
\]
$D_c^{s}$ is robotaxi choice demand \emph{before} any vehicle shortage is applied.
$D_c^{N,s}$ covers human-driven taxis, other modes, and giving up the trip. In this baseline model, $M_c$ is assumed to be fixed.

\subsection{Outcome variables}

Given $D_c^{s}$ and service capacity $\kappa_c q_c$, the following quantities
(both in trips) follow from the deployment outcome and do not need to be
estimated separately. Here $(x)^+=\max\{x,0\}$.
\[
O_c^{s}=(\kappa_c q_c-D_c^{s})^{+}\ \ \text{(unused capacity)},\qquad
U_c^{s}=(D_c^{s}-\kappa_c q_c)^{+}\ \ \text{(unserved demand)}.
\]

%%%%%%%%%%%%%%%%%%%%%%%%%%%%%%%%%%%%%%%%%%%%%%%%%%%%%%
\section{Notation}

{\small
\begin{tabularx}{\linewidth}{@{}>{\raggedright\arraybackslash}p{0.09\linewidth}>{\raggedright\arraybackslash}p{0.40\linewidth}Y@{}}
\toprule
Symbol & Meaning & Type and unit \\ \midrule
$C$ & Set of candidate cities & Input set; cities \\
$S$ & Set of demand scenarios & Input set; scenarios \\
$y_c$ & 1 if city $c$ is entered, else 0 & Decision variable; binary \\
$q_c$ & Number of vehicles deployed in city $c$ & Decision variable; vehicles (integer) \\
$\kappa_c$ & Trips one vehicle can serve during the period & Parameter; trips/vehicle per period \\
$Q$ & Total vehicles available for deployment & Parameter; vehicles \\
$\bar q_c$ & Upper bound on deployment in city $c$ & Parameter; vehicles \\
$F_c$ & Fixed cost of city entry and basic infrastructure & Parameter; currency per period \\
$K_c$ & Per-vehicle preparation, deployment, and holding cost & Parameter; currency/vehicle per period \\
$c_c^{+}$ & Additional loss per unit of unused capacity & Parameter; currency/unused trip \\
$c_c^{-}$ & Foregone contribution margin per unit of unserved demand & Parameter; currency/unserved trip \\
$M_c$ & Potential ride-hailing usage opportunities in city $c$ & Parameter; opportunities per period (data or assumption) \\
$p_c^{s}$ & Probability of choosing the company's robotaxi & Scenario input; dimensionless, $[0,1]$ \\
$D_c^{s}$ & Robotaxi choice demand, $M_c p_c^s$ & Derived input; trips per period \\
$D_c^{N,s}$ & Non-robotaxi demand, $M_c(1-p_c^s)$ & Derived input; trips per period \\
$O_c^{s}$ & Unused service capacity & Outcome variable; trips per period \\
$U_c^{s}$ & Unserved robotaxi demand & Outcome variable; trips per period \\
\bottomrule
\end{tabularx}
}

\textbf{Notes.}
\begin{itemize}[leftmargin=*,itemsep=1pt]
  \item Demand inputs are $M_c$ and $p_c^s$; \textcolor{blue}{in the baseline model $M_c$ is fixed and
        uncertainty enters only through $p_c^s$.}
  \item Outcome variables ($O_c^s$, $U_c^s$) are computed from demand and deployment,
        not predicted from separately collected data.
\end{itemize}


%%%%%%%%%%%%%%%%%%%%%%%%%%%%%%%%%%%%%%%%%%%%%%%%%%%%%%
\section{Key assumptions and open points}
\begin{itemize}[leftmargin=*,itemsep=1pt]
  \item A single evaluation period must be fixed; all parameters are expressed over it.
  \item $p_c^s$ must be specified relative to reference service conditions (e.g.\ fares, waiting times).
  \item $M_c$ is not directly observable; it is constructed from observed-usage patterns and assumptions.
  \item Scenarios are equally likely.
\end{itemize}

\end{document}